\documentclass[12pt]{exam}

% --- Packages ---
\usepackage[utf8]{inputenc}
\usepackage[T1]{fontenc}
\usepackage[french]{babel}
\usepackage{amsmath, amssymb}
\usepackage{geometry}
\usepackage{lmodern}

\geometry{margin=2cm}
\setlength{\parindent}{0pt}

% =====================
% Document
% =====================
\begin{document}

\begin{center}
    {\Large \textbf{Feuille d’exercices de retenue}}\\[2mm]
    {\large 2nde}\\[2mm]
    \textit{Équations}
\end{center}

\vspace{0.3cm}

\textbf{Consigne générale :}  
Résoudre les équations suivantes.

\vspace{0.4cm}

% ==================================================
\begin{questions}

% -----------------------------
\question \textbf{Équations simples}

Résoudre les équations suivantes :
\begin{parts}
\part $3x + 5 = 14$
\vspace{0.3cm}
\part $7x - 4 = 3x + 8$
\vspace{0.3cm}
\part $2(x - 3) = 10$
\end{parts}

\vspace{0.5cm}

% -----------------------------
\question \textbf{Équations avec fractions}

Résoudre les équations suivantes :
\begin{parts}
\part $\dfrac{x}{3} = 4$
\vspace{0.3cm}
\part $\dfrac{2x - 1}{5} = 3$
\vspace{0.3cm}
\part $\dfrac{3}{4}x + 2 = 5 - \dfrac{2}{3}x$
\end{parts}

\vspace{0.5cm}

% -----------------------------
\question \textbf{Équations avec fractions (mise au même dénominateur)}

Résoudre les équations suivantes :
\begin{parts}
\part $\dfrac{x}{3} = \dfrac{2}{5}$
\vspace{0.3cm}
\part $\dfrac{2x + 1}{4} = \dfrac{x - 3}{2}$
\end{parts}

\vspace{0.5cm}

% -----------------------------
\question \textbf{Équations produits}

Résoudre les équations suivantes :
\begin{parts}
\part $(x - 3)(x + 2) = 0$
\vspace{0.3cm}
\part $(2x - 1)(x - 4) = 0$
\vspace{0.3cm}
\part $x(x - 5) = 0$
\vspace{0.3cm}
\part $(x + 1)(3x - 6) = 0$
\end{parts}

\end{questions}

\vspace{0.5cm}

\begin{center}
    \textit{Fin de la feuille}
\end{center}

\end{document}
